% Part: lambda-calculus
% Chapter: introduction
% Section: basic-pr-lambda

\documentclass[../../../include/open-logic-section]{subfiles}

\begin{document}

\olfileid{lam}{int}{bas}
\olsection{મૂળભૂત આદિમ પુનરાવર્તી વિધેયો \usetoken{S}{lambda definable} છે}

\begin{lem}
વિધેયો $\Zero$, $\Succ$ અને $\Proj{n}{i}$ !!{lambda definable} છે.
\end{lem}

\begin{proof}
$\Zero$ એટલે માત્ર $\lambd[x][\lambd[y][y]]$.

ઉત્તરગામી વિધેય~$\Succ$ને
$\fn{Succ}(u) = \lambd[x][\lambd[y][x(uxy)]]$ વડે વ્યાખ્યાયિત
કરીએ છીએ. આ કેવી રીતે કામ કરે છે તે વિચારો. દરેક ચર્ચ અંક~$\num{n}$ને
પુનરાવર્તક માનો અને કોઈપણ વિધેય~$f$ લો. ત્યારે
$\fn{Succ}(\num{n},f)$ એવું વિધેય છે જે ઇનપુટ~$y$ પર, $y$થી
શરૂ કરીને~$f$ને $n$ વખત લાગુ કરે છે અને પછી વધુ એક વખત લાગુ કરે છે.

પ્રક્ષેપો માટે કંઈ વિશેષ કહેવાનું નથી:
$\fn{Proj}^n_i(x_0, \dots, x_{n-1}) = x_i$. બીજા શબ્દોમાં,
આપણા સંકેતન મુજબ $\fn{Proj}^n_i$ એ લૅમ્બડા પદ
$\lambd[x_0][\dots \lambd[x_{n-1}][x_i]]$ છે.
\end{proof}

\end{document}
